\section{Redes e filtros}

\subsection*{Convergência de seqüências}

\begin{definition}
    Seja \((X,d)\) e \((x_n)_{n\in \N}\subset X\) uma seqüência. Dizemos que \(x_n \to x \in X\) (\emph{converge}) sse \(\forall \epsilon >0, \ \exists n_0 \in \N : d(x_n, x) <\epsilon \) se \(n\geq n_0\).
\end{definition}

\begin{proposition}
    \textbf{12.2.} Sejam \(A \subset (X, d)\) e \(x\in X\). Então, \(x \in \overline{A} \ \leftrightharpoons \ \exists (x_n) \subset A : x_n \to x\). 
\end{proposition}

\begin{corollary}
    \textbf{12.3.} \(A\ce \subset (X,d) \ \leftrightharpoons \ \forall (x_n)\subset A: x_n \to x\) tem-se \(x\in A\). 
\end{corollary}

\begin{proposition}
    \textbf{12.4.} Uma função \(f : (X, d_X) \to (Y, d_Y)\) é contínua em \(a\in X\) \ \(\leftrightharpoons\) \ Pra cada \(x_n \to a \) (em \(X\)) tem-se que \(f(x_n) \to f(a)\) (em \(Y\)). \textcolor{gray}{(\(\Leftarrow\)) Contrapositiva}
\end{proposition}

\begin{definition}
    Dizemos que \(x_n \to x\) (em \((X, \tau)\)) sse \(\forall U \in \U_x, \ \exists n_0 \in \N : (x_n) \subset U \) se \(n \geq n_0\).
\end{definition}

\begin{definition}
    Chamamos de \emph{subseqüência} de \((x_n)_{n\in \N}\) qualquer \((x_{n_k})_{k\in \N} : (n_k)_{k\in \N}\) é uma seqüências estritamente crescente em \(\N\). 
\end{definition}

\subsection*{Convergência de redes}

\begin{definition}
    Um duplo \((\Lambda, \leq)\) é chamado de \emph{conjunto dirigido} se a relação verificar (D1) reflexividade; (D2) transitividade e (D3) \(\forall \lambda, \mu \in \Lambda, \ \exists \nu \in \Lambda : \lambda, \mu \leq \nu\).
\end{definition}

\begin{definition}
    Sejam \(\Lambda\) dirigido e \((X, \tau)\). Chamamos de \emph{rede} qualquer função \(\x : \Lambda \to X\), escrevemos \( (x(\lambda))_{\lambda\in \Lambda} = (x_\lambda)_{\lambda \in \Lambda} \subset X\). Dizemos que \(x_\lambda \to x \in X\) (\emph{converges}) sse \(\forall U \in \U_x: \exists \lambda_0 \in \Lambda : (x_\lambda) \subset U\) se \(\lambda \geq \lambda_0\).  
\end{definition}

\ejemplo{\centering Vizinhanças dirigidas
\tcblower
\begin{example}
    Dadas \(U, V \in \U_x\) defina \(U \leq V \ \leftrightharpoons \ V \subset U\). Assim definido \((\U_x, \leq) \)  é dirigido. Seja \((x_U)_{U\in \U_x} : x_U \in U \in \U_x\), é claro que \(x_U \to x\).  
\end{example}
}

\begin{proposition}
    \textbf{13.5.} Sejam \(A \subset (X, \tau)\) e \(x \in X\). Então, \(x \in \overline{A} \ \leftrightharpoons \ \exists (x_\lambda)\subset A : x_\lambda \to x\). \comentario{Prova as duas direções usando a \(\square \) \hyperlink{prop58}{5.8.}}
\end{proposition}

\hypertarget{prop136}{}
\begin{proposition}
    \textbf{13.6.} Uma função \(f : (X, \tau_X) \to (Y, \tau_Y)\) é contínua em \(a \in X \ \leftrightharpoons \ \) Pra cada \(x_\lambda \to a\) (em \(X\)) tem-se que \(f(x_\lambda) \to f(a)\) (em \(Y\)). \textcolor{gray}{(\(\Leftarrow\)) Contrapositiva}
\end{proposition}

\begin{proposition}
    \textbf{13.7.} Sejam \((\x_\lambda) \subset X := \prod X_i\) e \(\x \in X\). Então, \(\x_\lambda \to \x \in X\) \ \(\leftrightharpoons\) \ \(\pi_i(\x_\lambda) \to \pi_i(\x)\) (em \( X_i\)), pra cada \(i \in I\).
\end{proposition}

\demo{(\(\Rightarrow\)) \(\square\) \hyperlink{prop136}{13.6.}; (\(\Leftarrow\)) Seja \(U \in \U_\x\) (em \(X\)), então \(\exists V \in \U_\x : \x\in V \subset U\) e, 
\[V = \bigcap_{j\in J} \pi_j^{-1}(V_j) : J \subset I \text{ finito  e } V_j\ab \subset X_j. \]
Em particular, \(V_j \in \U_{\pi_j(\x)}\) (em \(X_j\)), logo \(\forall j \in J, \ \exists \lambda_j \in \Lambda : \pi(\x_\lambda) \subset V_j \) se \(\lambda \geq \lambda_j\). Finalmente, \(\exists \lambda_0 \in \Lambda : \lambda_0 \geq \lambda_j, \ \forall j \in J \) e \((\x_\lambda) \subset V \subset U\) se \(\lambda \geq \lambda_0\).% e \(\x_\lambda \to \x\). 
}

\begin{definition}
    Um ponto \(x \in X\) é \emph{ponto de acumulação} de \((x_\lambda) \subset X\) se dados \(U \in \U_x\) e \(\lambda_0 \in \Lambda, \ \exists \lambda \geq \lambda_0\) tal que \( x_\lambda \in U\).
\end{definition}

\begin{definition}
    Chamamos de \emph{subrede} de \(\x : \Lambda \to X\) qualquer rede da forma \(\x \circ \phi : M \to X\), denotada \((x_{\phi(\mu)})_{\mu \in M}\), onde \(M \) é dirigido e \(\phi:M \to \Lambda\) verifica:
    \vspace{-0.2cm}
    \begin{itemize}[left = 1.8cm]
        \item[(Sr1)] \(\mu_1 \leq \mu_2 \Rightarrow \phi(\mu_1) \leq \phi(\mu_2) \ \leftarrow\) \ Crescente 
        \item[(Sr2)] \(\forall  \lambda \in \Lambda, \ \exists \mu \in M : \phi(\mu )\geq \lambda\ \leftarrow\) \ \emph{Cofinal}
    \end{itemize}
\end{definition}

\begin{proposition}
    \textbf{13.10.} Sejam \((x_\lambda) \subset X\) e \(x \in X\). Então, \(x\) é ponto de acumulação de \((x_\lambda)\) \ \(\leftrightharpoons\) \ \(\exists (x_{\phi(\mu)})\) subrede de \((x_\lambda)\) tal que \(  x_{\phi(\mu)} \to x\). 
\end{proposition}

\demo{(\(\Rightarrow\)) Defina \(M := \{(\lambda, U) \in \Lambda \times \U_x : x_\lambda\in U\}\), \(\phi : M \to \Lambda \) tal que \((\lambda, U) \mapsto \lambda\) e mostra que \((x_{\phi(\mu)})\) é subrede de \((x_\lambda) : x_{\phi(\mu)} \to x\); (\(\Leftarrow \)) Segue das definições.}

\begin{definition}
    Uma rede \((x_\lambda) \subset X\) é dita \emph{universal} se dado \(A \subset X, \ \exists \lambda_0 \in \Lambda \) tal que \(\{x_\lambda : \lambda \geq \lambda_0\} \subset A\) ou \(\{x_\lambda : \lambda \geq \lambda_0\} \subset X \setminus A\).
\end{definition}

\begin{proposition}
    \textbf{13.12.} Sejam \((x_\lambda) \subset X\) universal e \(x\in X\). Se \(x\) é ponto de acumulação de \((x_\lambda)\), então \(x_\lambda \to x\). \comentario{Segue direto das definições}
\end{proposition}

%\begin{definition}
%    Uma conjunto \((X, \leq)\) é \emph{parcialmente ordenado} se a \emph{relação} \(\leq\) for (P$1$) reflexiva; (P$2$) antisimétrica e (P$3$) transitiva. %É \emph{total} se +(P$4$) \(\forall x,y \in X, \ x\leq y\) ou \(y\leq x\).
%\end{definition}
%%
%\begin{example}
%    \((\mathcal{P}(X), \subseteq)\) e \((\R, \leq)\) são parcialmente ordenados. 
%\end{example}
%
%\begin{definition}
%    Sejam \(X\) parcialmente ordenado e \(A \subset X\), 
%    \begin{itemize}[label = \(\star\)]
%        \item Se \(\exists a_0 \in A : a_0 \leq a, \ \forall a \in A  \ \leftarrow \) \ \emph{Elemento mínimo} (ou máximo resp.)
%        \item Se \(\exists a_0 \in A : a \leq a_0 \Rightarrow a = a_0, \ \forall a \in A  \ \leftarrow \) \ \emph{Elemento mínimo}
%    \end{itemize}
%\end{definition}

\hypertarget{lemma145}{}
\begin{lemma}
    \textbf{14.5. (\emph{Zorn})} Seja \(X\neq \varnothing\) parcialmente ordenado tal que cada cadeia em \(X\) é limitada superiormente. Então \(X\) possui pelo menos um elemento maximal.
\end{lemma}

\hypertarget{thm146}{}
\begin{theorem}
    \textbf{14.6. (\emph{Zermelo})} Todo conjunto \(X\neq \varnothing\) pode ser bem ordenado.
\end{theorem}

\begin{theorem}
    \textbf{14.7.} Axioma da escolha \ \(\leftrightharpoons\) \ {\scriptsize $\square$}\ \hyperlink{lemma145}{14.5. (\emph{Zorn})} \ \(\leftrightharpoons\) \ \(\blacksquare\) \ \hyperlink{thm146}{14.6. (\emph{Zermelo})}. \comentario{Não suporta resumo, olha se queser \cite[pag. 39]{mujica2005}}
\end{theorem}

\subsection*{Convergência de filtros}

\begin{definition}
    Sejam \(X \neq \varnothing\). Dizemos que uma familia não vazia \(\neq \mathcal{F}\subset \mathcal{P}(X)\) é um \emph{filtro} em \(X\) se verificar as seguintes condições: 
    \begin{enumerate}[left = 3cm]
        \item[(F$1$)] \(A \neq \varnothing, \ \forall A \in \mathcal{F}\). \hspace{1cm} (F$2$) \(A, B \in \mathcal{F} \Rightarrow A\cap B \in \mathcal{F}\).
        \item[]\hspace{1cm}(F$3$) \(A \in \mathcal{F}\) e \(A \subset B \subset X \Rightarrow B \in \mathcal{F}\). 
    \end{enumerate}
\end{definition}